\documentclass[12pt]{amsart}
\usepackage{amsmath,amssymb,amsthm}
\usepackage{hyperref}

\newtheorem{theorem}{Theorem}[section]
\newtheorem{proposition}[theorem]{Proposition}
\newtheorem{lemma}[theorem]{Lemma}
\newtheorem{definition}[theorem]{Definition}
\newtheorem{example}[theorem]{Example}
\newtheorem{remark}[theorem]{Remark}

\newcommand{\Ind}{\operatorname{Ind}}
\newcommand{\Res}{\operatorname{Res}}
\newcommand{\Sym}{\operatorname{Sym}}
\newcommand{\sgn}{\operatorname{sgn}}
\newcommand{\triv}{\operatorname{triv}}
\newcommand{\refl}{\operatorname{refl}}
\newcommand{\Irr}{\operatorname{Irr}}
\newcommand{\JInd}{j\text{-}\!\operatorname{Ind}}
\newcommand{\TInd}{\operatorname{TGInd}}
\newcommand{\FF}{\mathbb{F}}
\newcommand{\ZZ}{\mathbb{Z}}
\newcommand{\QQ}{\mathbb{Q}}
\newcommand{\CC}{\mathbb{C}}
\newcommand{\bfG}{\mathbf{G}}
\newcommand{\bfL}{\mathbf{L}}
\newcommand{\bfT}{\mathbf{T}}
\newcommand{\bfW}{\mathbf{W}}
\newcommand{\cF}{\mathcal{F}}
\newcommand{\cO}{\mathcal{O}}

\title{Generic Degrees and Special Representations of Weyl Groups}
\author{Notes on the Atlas Software Implementation}
\date{\today}

\begin{document}

\maketitle
\tableofcontents

\section{Introduction}

Let $\bfG$ be a connected reductive algebraic group over an algebraically
closed field, with Weyl group $W$. The irreducible complex representations of
$W$ play a central role in the representation theory of $\bfG$, both over
finite fields (where they parametrize unipotent characters via Green functions)
and in the study of primitive ideals in the universal enveloping algebra of the
Lie algebra of $\bfG$.

Two numerical invariants are attached to each irreducible $W$-representation
$E$:
\begin{itemize}
  \item the \emph{degree} (or \emph{$b$-value}) $b(E)$, which is an integer
        measuring the first occurrence of $E$ in the symmetric powers of the
        reflection representation;
  \item the \emph{generic degree} (or \emph{$a$-value}) $a(E)$, introduced by
        Lusztig, which controls the pole order of $E$ in the generic Hecke
        algebra.
\end{itemize}
The inequality $a(E) \leq b(E)$ always holds, and equality $a(E) = b(E)$
characterizes a distinguished class of representations called \emph{special
representations}.  These notions organize the irreducible $W$-representations
into \emph{families}, each containing exactly one special representation.

This document gives an account of these definitions and the algorithms
implemented in the Atlas software to compute them.  The primary references are
Lusztig~\cite{Lusztig:orange}, Carter~\cite{Carter:finite}, and
Geck~\cite{Geck:chevie}.

\section{Weyl Groups and Their Representations}

\subsection{Setup}

Let $\bfG$ be a connected reductive group with maximal torus $\bfT$, root
system $\Phi \subset X^*(\bfT)$, and Weyl group $W = N_{\bfG}(\bfT)/\bfT$.
Fix a set of positive roots $\Phi^+$ with simple roots $\alpha_1,\ldots,
\alpha_r$.  For each $\alpha \in \Phi$ let $s_\alpha \in W$ denote the
corresponding reflection.

The group $W$ is a finite Coxeter group with generators $\{s_i = s_{\alpha_i}
\mid 1 \leq i \leq r\}$.  We write $\Irr(W)$ for the set of isomorphism
classes of irreducible complex representations of $W$.

\subsection{The reflection representation}

The \emph{reflection representation} $\refl$ of $W$ is the natural action of
$W$ on $\mathfrak{h}^* = X^*(\bfT) \otimes_\ZZ \QQ$, which is faithful and
irreducible for simple $W$.

\section{The Degree ($b$-Value)}

\subsection{Definition}

The \emph{module of coinvariants} of $W$ is $\CC[V] / (C[V]^W_+)$ where
$V = \mathfrak{h}^*_\CC$ and $(C[V]^W_+)$ is the ideal generated by
$W$-invariant polynomials of positive degree.  This module is graded and affords
a graded $W$-module structure; its Poincar\'e series encodes the
\emph{fake degrees}.

More concretely, for each $E \in \Irr(W)$ the \emph{fake degree polynomial}
$f_E(q) \in \ZZ[q]$ is defined by the multiplicity of $E$ in the graded module
of coinvariants:
\[
  f_E(q) = \sum_{k=0}^{|\Phi^+|} \langle E,\, \Sym^k(\refl) \rangle_W \, q^k.
\]
(Here $\langle \cdot, \cdot \rangle_W$ denotes the inner product of characters
of $W$.)

\begin{definition}
  The \emph{degree} (or \emph{$b$-value}) of $E \in \Irr(W)$ is
  \[
    b(E) = \text{lowest degree of } q \text{ in } f_E(q)
    = \min\{ k \geq 0 \mid \langle E,\, \Sym^k(\refl) \rangle_W \neq 0 \}.
  \]
\end{definition}

\subsection{Implementation in Atlas}

In the Atlas software (see \texttt{character\_tables.at}), the character table
of $W$ is stored as a \texttt{CharacterTable} structure.  The field
\texttt{deg\_spec} stores pairs $(b(E), \text{special}(E))$ for each irrep,
computed at construction time.  The degree $b(E)$ is found by searching for the
smallest $k$ such that $\langle E, \Sym^k(\refl) \rangle \neq 0$:
\begin{verbatim}
  first(limit, (int k) bool: !=Wct.inner(char, Wct.sym_power_refl(k)))
\end{verbatim}
The function \texttt{degree(ct)(i)} then returns the $b$-value of irrep $i$.

\section{Special Representations and Families}

\subsection{The Springer Correspondence}

The Springer correspondence \cite{Springer:1978} gives a bijection between the
set of nilpotent orbits $\cO$ of $\bfG$ (or more precisely, pairs $(\cO, \phi)$
where $\phi$ is an irreducible character of the component group of the
centralizer of a point of $\cO$) and a subset of the set of irreducible
$W$-representations.

For classical groups and in characteristic $0$ or large characteristic, this
gives a bijection between nilpotent orbits and a subset of $\Irr(W)$.
Explicitly:

\begin{itemize}
  \item \textbf{Type $A_{n-1}$} ($W = S_n$): Irreps are indexed by partitions
        of $n$; nilpotent orbits (Jordan types) are also indexed by partitions
        of $n$.  The Springer correspondence sends each orbit directly to the
        corresponding irrep.

  \item \textbf{Type $B_n$ / $C_n$}: Irreps are indexed by bipartitions
        $(\lambda, \mu)$ (ordered pairs of partitions with $|\lambda| +
        |\mu| = n$); nilpotent orbits are indexed by partitions of $2n$ (for
        $C_n$, those with each odd part occurring an even number of times;
        for $B_n$, each even part occurring an even number of times). The
        Springer map (via the core-quotient construction) sends each orbit
        to a bipartition.

  \item \textbf{Type $D_n$}: Similar to type $B$/$C$, with a complication
        for ``split'' classes and representations.

  \item \textbf{Exceptional types}: The Springer correspondence is given by
        explicit tables (see \cite{Carter:finite} and the Atlas files
        \texttt{springer\_table\_E6.at}, etc.).
\end{itemize}

\subsection{Special Orbits}

A nilpotent orbit $\cO$ is \emph{special} (in the sense of Lusztig) if it
appears in the image of the Springer correspondence on the principal series,
i.e., if $(\cO, \triv)$ is in the image.  Equivalently (for simply connected
$\bfG$ in good characteristic), $\cO$ is special if and only if the special
closure of $\cO$ equals $\cO$ itself.

For type $B_n$/$C_n$, a partition (nilpotent orbit) is special if and only if
its parts satisfy the appropriate parity conditions:
\begin{itemize}
  \item Type $C_n$: each odd part occurs an even number of times.
  \item Type $B_n$: each even part occurs an even number of times.
\end{itemize}

\subsection{Special Representations}

Via the Springer correspondence, each special orbit $\cO$ corresponds to an
irreducible $W$-representation $E_\cO = \text{Springer}(\cO)$.  The
representation $E_\cO$ is called a \emph{special representation} of $W$.

Equivalently, $E \in \Irr(W)$ is special if and only if $a(E) = b(E)$, i.e.,
the generic degree equals the degree.

\begin{definition}
  An irreducible representation $E \in \Irr(W)$ is \emph{special} if
  $a(E) = b(E)$.
\end{definition}

In the Atlas software, the function $\texttt{to\_special}: \Irr(W) \to \Irr(W)$
sends each irrep to its associated special representation (defined via the
Springer map and Lusztig's $\mathbf{a}$-function theory).  The identity
\[
  \texttt{is\_special}(E) \iff \texttt{to\_special}(E) = E
  \iff \texttt{special}(E) = E
\]
holds.  For exceptional types, \texttt{to\_special} is stored in explicit
lookup tables (e.g., \texttt{to\_special\_E8} in
\texttt{character\_table\_E8.at}):
\begin{verbatim}
  set to_special_G2 = (int i) int:
     case i in 0, 5, 5, 3, 5, 5 else error("Wrong G2 irrep number") esac
\end{verbatim}
For classical types, \texttt{to\_special} is computed by the
\texttt{make\_special} function on bipartitions or symbols (in
\texttt{combinatorics.at}).

\subsection{Families}

\begin{definition}
  Two irreducible representations $E, E' \in \Irr(W)$ are in the same
  \emph{family} if $\texttt{to\_special}(E) = \texttt{to\_special}(E')$,
  i.e., if they share the same associated special representation.
\end{definition}

This defines a partition of $\Irr(W)$ into families.  Each family contains
exactly one special representation.  In the Atlas software
(\texttt{families.at}):
\begin{verbatim}
  set families(CharacterTable ct)=[[int]]:
  let v = for i:ct.n_classes do [int]: [] od
  in for j:ct.n_classes do v[ct.special(j)] #:= j od
  ;  ...sort_special_first(ct,a)...
\end{verbatim}
The function \texttt{families} returns a list of families, each sorted with the
special representation first.

\section{Generic Degree ($a$-Value)}

\subsection{Definition via Families}

\begin{definition}
  The \emph{generic degree} (or \emph{$a$-value}) of $E \in \Irr(W)$ is
  \[
    a(E) = b(E_s)
  \]
  where $E_s$ is the unique special representation in the family of $E$.
\end{definition}

Thus: $a(E) \leq b(E)$ with equality if and only if $E$ is itself special.  All
representations in the same family share the same $a$-value.

In the Atlas software (\texttt{character\_tables.at}, line 148):
\begin{verbatim}
  set generic_degree (CharacterTable ct) = (int i)int: ct.degree(ct.special(i))
\end{verbatim}

\subsection{Definition via the Generic Hecke Algebra}

Lusztig's original definition of the $a$-function is more algebraic.  Let
$\mathcal{H}$ be the generic Hecke algebra of $W$ over $\ZZ[v, v^{-1}]$, with
standard basis $\{T_w\}_{w \in W}$.  Each irreducible $W$-representation $E$
gives rise to an irreducible $\mathcal{H}$-module $E_v$ over the fraction field.
The \emph{generic degree} $a(E)$ is defined so that $v^{a(E)} \cdot
c_{\text{leading}}$ is a unit in the $\ZZ[v]$-form of the module, where
$c_{\text{leading}}$ comes from the structure constants in the Kazhdan--Lusztig
basis $\{C_w\}$.  More precisely, $a(E)$ is the minimum non-negative integer
such that $v^{a(E)} \langle E_v, C_w \rangle$ lies in $\ZZ[v]$ for all $w \in
W$ \cite[Chapter 5]{Lusztig:orange}.

\subsection{Connection with Kazhdan--Lusztig Cells}

A fundamental theorem of Lusztig \cite{Lusztig:orange} connects generic degrees
with the Kazhdan--Lusztig cell theory:

\begin{theorem}[Lusztig]
  \begin{enumerate}
    \item Each two-sided Kazhdan--Lusztig cell $\mathbf{c}$ of $W$ contains
          exactly one special representation $E_s$.
    \item All irreps $E$ appearing in the cell character of $\mathbf{c}$ satisfy
          $a(E) = b(E_s)$.
    \item The special representation $E_s$ in $\mathbf{c}$ is the unique irrep
          of minimum $b$-value among those appearing in the cell character of
          any left cell contained in $\mathbf{c}$.
  \end{enumerate}
\end{theorem}

The correspondence $\mathbf{c} \leftrightarrow E_s$ gives a bijection between
the set of two-sided cells of $W$ and the set of special representations; thus
the families of $\Irr(W)$ are in bijection with two-sided cells.

\subsection{Connection with Green Functions and Unipotent Characters}

When $\bfG = \bfG(\FF_q)$ is a finite group of Lie type, the generic degree
$a(E)$ appears in the pole order of the generic degree polynomial of the
unipotent character associated to $E$ via the Deligne--Lusztig theory.  The
Green functions $Q^{\bfG}_\bfL$ encode the Springer correspondence and the
$a$-values control the leading behavior in $q$ of matrix coefficients.

\section{Computing Generic Degrees: Two Algorithms}

The Atlas software implements two approaches to computing generic degrees.

\subsection{Cell-Based Algorithm}

The first algorithm computes generic degrees from Kazhdan--Lusztig cells
(implemented in \texttt{generic\_degrees.at}):

\begin{enumerate}
  \item Compute the Kazhdan--Lusztig left cells of $W$ (via the built-in
        \texttt{W\_cells\_of} function, which uses the KL polynomials of a
        block of the representation-theoretic category $\mathcal{O}$).
  \item For each left cell $C$, compute the \emph{cell character}: the
        character of $W$ acting on the cell module,
        $\chi_C = \sum_{w \in C} (-1)^{\ell(w)} \chi_w$ (in appropriate sense).
  \item Decompose $\chi_C$ into irreducibles: $\chi_C = \sum_E m_E(C) \cdot E$.
  \item Identify the \emph{special irreducible} in $C$: the unique irrep of
        smallest $b$-value among those with $m_E(C) \neq 0$:
        \begin{verbatim}
  set special_irreducible (CharacterTable ct, [int] multiplicities) = int:
     let summands = smallest_degree_summands(ct, multiplicities)
  in case #summands
     in error("zero character, no special")
     , summands[0]
     else error("non-unique special representation")
     esac
        \end{verbatim}
  \item The generic degree of every irrep $E$ appearing in $\chi_C$ is
        $a(E) = b(E_s)$ where $E_s$ is this special irreducible.
\end{enumerate}

This method is efficient when the cells can be computed (which requires
computing KL polynomials of a block), but may be slow or infeasible for very
large rank groups (e.g., $E_8$) due to the size of the blocks.

\subsection{Geck's Inductive Algorithm}

The second algorithm, due to Geck \cite{Geck:chevie}, computes the $a$-function
inductively on the rank (implemented in \texttt{geck\_generic.at}).

\paragraph{Setup.} For a Weyl group $W$ with root system $\Phi$, let
$\mathcal{L}$ denote the set of proper parabolic subgroups (Levi factors of
proper parabolic subgroups), i.e., the Weyl groups $W_J = W(\Phi_J)$ for
$J \subsetneq \{1, \ldots, r\}$.  The inductive hypothesis is that $a(E')$ is
known for all proper subgroups $W_J$ and all $E' \in \Irr(W_J)$.

\paragraph{Key quantity.} For $E \in \Irr(W)$, define
\[
  a'(E) = \max \bigl\{ a(E') \mid L \in \mathcal{L},\, E' \in \Irr(W_L),\,
                        \langle \Ind_{W_L}^W(E'), E \rangle \neq 0 \bigr\},
\]
i.e., $a'(E)$ is the maximum $a$-value over all proper Levi subgroup
representations that appear in the induced representation.  In Atlas
(\texttt{geck\_generic.at}):
\begin{verbatim}
  set a_prime(RootDatum G, [CharacterTable] tables,
              int table_index, int character_index) = int:
  max( for M in max_Levi_subgroups(G) do
         for char_M@t in ct_M.characters do
           if ind[character_index]>0 then [ct_M.generic_degree(t)] else [] fi
         od
       od )
\end{verbatim}

\paragraph{The $\omega_L$ correction.} For $E \in \Irr(W)$, define a
correction term involving the character values at reflections.  Following Geck
\cite[Definition 4.1]{Geck:chevie}:
\[
  \omega_L(E) = \frac{1}{\dim E} \sum_{C_s}
    |C_s| \cdot \chi_E(s)
\]
where the sum runs over conjugacy classes $C_s$ of reflections $s \in W$, and
$\chi_E$ denotes the character of $E$.  (Here $|C_s|$ is the size of the
class.)

\paragraph{The formula.} Geck's formula is:
\begin{equation}
  \label{eq:geck}
  a(E) = \max \bigl( a'(E),\; a'(E \otimes \sgn) - \omega_L(E) \bigr)
\end{equation}
where $\sgn$ is the sign character of $W$.

In Atlas (\texttt{geck\_generic.at}):
\begin{verbatim}
  let A = a_prime(G, tables, table_index, character_index)
  ,   B = a_prime(G, tables, table_index,
                  table.tensor_sign_index(character_index))
  ,   omega_L = omega_L(table, character_index, ...)
  then max = max(A, B - omega_L)
  in table.set_generic_degree(character_index, d, true)
\end{verbatim}

\paragraph{Inductive structure.} The algorithm builds a list of character
tables for all Levi subgroups (from small to large), filling in generic degrees
inductively.  The function \texttt{all\_levi\_tables} computes tables for all
Levi subgroups, ordered so that proper Levis of any given group appear later in
the list:
\begin{verbatim}
  set set_generic_degrees_geck(RootDatum G) = CharacterTable:
    update_generic_degrees_geck(all_levi_tables(G))
\end{verbatim}

\subsection{Mixed Strategy}

In practice, Atlas uses a combined strategy (\texttt{set\_generic\_degrees\_3}
in \texttt{geck\_generic.at}):
\begin{enumerate}
  \item First attempt to compute generic degrees using the cell algorithm for
        all Levi subgroups where cells are feasible.
  \item Then use the sign symmetry: since $a(E \otimes \sgn) = a(E) +
        \omega_L(E)$, knowing $a(E)$ immediately gives $a(E \otimes \sgn)$
        (function \texttt{update\_generic\_degrees\_using\_sign}).
  \item Finally, apply Geck's formula to any remaining irreps.
\end{enumerate}
\begin{verbatim}
  set set_generic_degrees_3 (RootDatum G) = [CharacterTable]:
  ( let tables = all_levi_tables_plus(G)
    in update_generic_degrees_using_sign(tables)
    ;  update_generic_degrees_geck(tables)
    ;  tables
  )
\end{verbatim}

\section{Algorithm for the Map to the Special Representation}

Given $E \in \Irr(W)$, a separate algorithm (in \texttt{special\_rep.at})
computes the map $E \mapsto E_s$ (the special representation in the same
family) without going through families.  This relies on the following theorem:

\begin{theorem}[Lusztig, {\cite{Lusztig:orange}}]
  \label{thm:to_special}
  \begin{enumerate}
    \item[\textup{(a)}] If $L$ is a Levi subgroup of $G$, and $C_L$ is a
          left cell of $W_L$, then $\TInd(C_L)$ is a left cell of $W$.
    \item[\textup{(b)}] If $C$ is a left cell of $W$, so is $C \otimes \sgn$.
    \item[\textup{(c)}] If $W \neq 1$, then for any left cell $C$ of $W$ there
          is a proper Levi subgroup $L$ and a left cell $C_L$ of $W_L$ such that
          either $\TInd(C_L) = C$, or $\TInd(C_L \otimes \sgn) = C \otimes \sgn$.
  \end{enumerate}
\end{theorem}

Here $\TInd$ denotes truncated (or $j$-)induction, which preserves cells and
families.  The algorithm to compute $E \mapsto E_s$ is:
\begin{enumerate}
  \item If $E$ is already special, return $E$.
  \item If $W$ is abelian, return the trivial representation.
  \item Use part (c) of the theorem: find a proper Levi $L$ and $E_L \in
        \Irr(W_L)$ such that $E$ appears in $j\text{-}\!\Ind_{W_L}^W(E_L)$
        with $a(E_L) = a(E)$.  By induction, compute $(E_L)_s$.  Then
        $E_s = j\text{-}\!\Ind_{W_L}^W((E_L)_s)$.
  \item If step 3 fails for $E$, replace $E$ with $E \otimes \sgn$, which is
        guaranteed to succeed; then apply the duality involution on special
        orbits to recover $E_s$.
\end{enumerate}

For exceptional types, the map $E \mapsto E_s$ is given by explicit tables
(e.g., in \texttt{character\_table\_G.at}, \texttt{character\_table\_E8.at},
etc.), while for classical types it is computed via the combinatorics of
bipartitions or symbols.

\section{Display Functions in Atlas}

The Atlas software provides several convenience functions for inspecting these
invariants:

\begin{itemize}
  \item \texttt{show(ct)}: Displays a table of all irreducible representations
        with their dimensions, the index of their special representation
        (marked with \texttt{*} if self-special), the $b$-value \texttt{deg},
        and the generic degree \texttt{gdeg}.
  \item \texttt{show\_degrees(ct)}: Alternate display comparing $a(i)$ and
        $b(i)$, with stars marking special representations.
  \item \texttt{families(ct)}: Returns the partition of $\Irr(W)$ into
        families, with each family sorted so the special representation appears
        first.
  \item \texttt{generic\_degree(ct)(i)}: Returns $a(E_i)$ for the $i$-th
        irrep.
\end{itemize}

\section{Example: Type $G_2$}

The Weyl group $W(G_2)$ has six irreducible representations, labeled
(in the Atlas conventions) $0, 1, \ldots, 5$:

\medskip
\begin{center}
\begin{tabular}{clccc}
\hline
$i$ & Name & $\dim$ & special($i$) & $a(E_i) = b(E_{\text{spec}})$ \\
\hline
0 & trivial              & 1 & 0\,(*) & $b(0)$ \\
1 & short sign           & 1 & 5      & $b(5)$ \\
2 & long sign            & 1 & 5      & $b(5)$ \\
3 & full sign            & 1 & 3\,(*) & $b(3)$ \\
4 & triangular reflection & 2 & 5     & $b(5)$ \\
5 & hexagonal reflection  & 2 & 5\,(*) & $b(5)$ \\
\hline
\end{tabular}
\end{center}
\medskip

The \texttt{to\_special\_G2} function (from \texttt{character\_table\_G.at}) is:
\begin{verbatim}
  set to_special_G2 = (int i) int:
     case i in 0, 5, 5, 3, 5, 5 else error("Wrong G2 irrep number") esac
\end{verbatim}
encoding the map $0\mapsto 0$, $1\mapsto 5$, $2\mapsto 5$, $3\mapsto 3$,
$4\mapsto 5$, $5\mapsto 5$.

The special representations are those marked (*): $0$ (trivial), $3$ (full
sign), and $5$ (hexagonal reflection).  The families are:
\begin{itemize}
  \item $\{0\}$, with special rep $0$;
  \item $\{3\}$, with special rep $3$;
  \item $\{1, 2, 4, 5\}$, with special rep $5$.
\end{itemize}
All four representations in the third family share the same generic degree
$a = b(5)$.

\begin{remark}
  In type $G_2$, the sign character splits into two distinct one-dimensional
  representations (short sign and long sign) reflecting the two root lengths.
  Both, together with the triangular reflection, lie in the family of the
  hexagonal reflection (the unique 2-dimensional special representation with
  the lower $b$-value), illustrating that non-singleton families can contain
  both one- and two-dimensional representations.
\end{remark}

\section{Summary}

The following diagram summarizes the relationships:

\[
\begin{array}{ccc}
  \text{nilpotent orbits of } G & \overset{\text{Springer}}{\longleftrightarrow} &
  \text{special reps of } W \\
  \updownarrow\text{\scriptsize special closure} & & \updownarrow\text{\scriptsize family} \\
  \text{special orbits} & \overset{\text{bijection}}{\longleftrightarrow} &
  \text{two-sided KL cells}
\end{array}
\]

For each $E \in \Irr(W)$:
\begin{align*}
  b(E) &= \min\{k \mid \langle E, \Sym^k(\refl)\rangle \neq 0\} \quad
           \text{(degree, } b\text{-value, fake degree exponent)} \\
  E_s  &= \text{unique special rep in family of } E
           \quad \text{(via \texttt{to\_special})} \\
  a(E) &= b(E_s) \quad \text{(generic degree, } a\text{-value)}
\end{align*}

The Atlas software computes $a(E)$ by either:
\begin{itemize}
  \item the \textbf{cell algorithm}: computing KL cells and finding the minimum-$b$
        irrep in each cell; or
  \item \textbf{Geck's inductive formula}
        $a(E) = \max(a'(E),\, a'(E \otimes \sgn) - \omega_L(E))$,
        proceeding inductively on Levi subgroups.
\end{itemize}
In practice a mixed strategy is used, applying the cell algorithm where
feasible and the Geck formula for the remainder.

\begin{thebibliography}{99}

\bibitem{Carter:finite}
R.~W. Carter,
\textit{Finite Groups of Lie Type: Conjugacy Classes and Complex Characters},
Wiley Classics Library, John Wiley \& Sons, 1985.

\bibitem{Geck:chevie}
M.~Geck,
\textit{Some applications of CHEVIE to the theory of algebraic groups},
Canad.\ J.\ Math.\ \textbf{61} (2009), 1239--1261.
(Cf.\ also: M.~Geck, Some applications of Chevie to the theory of algebraic
groups, preprint.)

\bibitem{KL:1979}
D.~Kazhdan and G.~Lusztig,
\textit{Representations of Coxeter groups and Hecke algebras},
Invent.\ Math.\ \textbf{53} (1979), 165--184.

\bibitem{Lusztig:orange}
G.~Lusztig,
\textit{Characters of Reductive Groups over a Finite Field},
Annals of Mathematics Studies \textbf{107},
Princeton University Press, 1984.

\bibitem{Lusztig:families}
G.~Lusztig,
\textit{A class of irreducible representations of a Weyl group},
Indag.\ Math.\ \textbf{41} (1979), 323--335;
\textit{A class of irreducible representations of a Weyl group, II},
Indag.\ Math.\ \textbf{44} (1982), 219--226.

\bibitem{Springer:1978}
T.~A. Springer,
\textit{A construction of representations of Weyl groups},
Invent.\ Math.\ \textbf{44} (1978), 279--293.

\end{thebibliography}

\end{document}
